\documentclass[12pt,a4paper]{article}

\usepackage[T1]{fontenc}
\usepackage[utf8]{inputenc}
\usepackage[catalan]{babel}
\usepackage[margin=2.5cm]{geometry}
\usepackage{lmodern}
\usepackage[nopatch=footnote]{microtype}
\usepackage{booktabs}
\usepackage{array}
\usepackage{tabularx}
\usepackage[table]{xcolor}
\usepackage[colorlinks=true,linkcolor=blue,urlcolor=blue]{hyperref}
\usepackage{fancyhdr}
\usepackage{setspace}

\definecolor{upcblue}{RGB}{0,51,153}

\newcommand{\InstructorName}{Oriol Farrés}
\newcommand{\CourseLevel}{1r de Batxillerat}
\newcommand{\AcademicYear}{2025--2026}
\newcommand{\DeadlineDate}{diumenge 3 de maig de 2026, a les 23:59}

\newcolumntype{Y}{>{\raggedright\arraybackslash}X}

\setstretch{1.12}
\setlength{\parskip}{0.45em}
\setlength{\parindent}{0pt}
\renewcommand{\arraystretch}{1.2}

\pagestyle{fancy}
\fancyhf{}
\fancyhead[L]{Rúbrica Lliurable 1}
\fancyhead[R]{\InstructorName}
\fancyfoot[C]{\thepage}
\setlength{\headheight}{14.5pt}
\renewcommand{\headrulewidth}{0.4pt}
\renewcommand{\footrulewidth}{0pt}

\begin{document}

\hypersetup{pageanchor=false}
\begin{titlepage}
  \thispagestyle{empty}
  \color{upcblue}
  \rule{\textwidth}{1.2pt}
  \color{black}

  \vspace{1.2cm}

  \begin{center}
    {\Large Rúbrica\par}
    \vspace{0.3cm}
    {\Huge\bfseries Lliurable 1 --- Document LaTeX\par}
    \vspace{0.35cm}
    {\large Criteris d'avaluació del document científic en LaTeX\par}
  \end{center}

  \vspace{0.6cm}

  \color{upcblue}
  \rule{\textwidth}{0.6pt}
  \color{black}

  \vspace{1.6cm}

  \begin{center}
    \begin{tabular}{>{\bfseries}p{4.2cm}p{9.2cm}}
      Instructor & \InstructorName \\
      Nivell & \CourseLevel \\
      Assignatura & Programació \\
      Curs acadèmic & \AcademicYear \\
      Data límit & diumenge 3 de maig de 2026 a les 23:59 \\
    \end{tabular}
  \end{center}

  \vfill

  \begin{center}
    \fbox{\parbox{0.9\textwidth}{
      \centering
      Aquest document defineix els criteris i nivells d'assoliment per avaluar el Lliurable 1.
      La nota màxima és 10 punts.
    }}
  \end{center}
\end{titlepage}
\hypersetup{pageanchor=true}

\section{Com s'utilitza aquesta rúbrica}

Aquest document estableix l'avaluació del \textbf{Lliurable 1}, que representa el \textbf{20\% de la nota final} de l'assignatura. Cada criteri s'avalua en un dels quatre nivells d'assoliment de la Taula~\ref{tab:nivells-rubrica}. El nivell \emph{Excel·lent} atorga la puntuació completa del criteri, mentre que la resta de nivells concedeixen una part proporcional del seu pes. Els punts extres s'atorguen manualment per part del professor només en treballs clarament excepcionals.

\begin{table}[htbp]
  \centering
  \caption{Nivells d'assoliment de la rúbrica}
  \label{tab:nivells-rubrica}
  {\small
  \begin{tabularx}{\textwidth}{>{\raggedright\arraybackslash}p{2.6cm}YY}
    \toprule
    \rowcolor{upcblue!15}
    \textbf{Nivell} & \textbf{Descriptor} & \textbf{Punts (relatiu al pes)} \\
    \midrule
    Excel·lent & Compleix tots els requisits amb claredat i precisió & 100\% del pes \\
    Notable & Compleix els requisits principals amb alguna mancança menor & 75\% del pes \\
    Aprovat & Compleix parcialment els requisits; errors o omissions notables & 50\% del pes \\
    Suspès & No compleix els requisits mínims o element absent & 25\% del pes \\
    \bottomrule
  \end{tabularx}
  }
\end{table}

\section{Bloc A --- Front Matter (2 punts)}

La Taula~\ref{tab:bloc-a} recull els criteris del bloc inicial del document.

\begin{table}[htbp]
  \centering
  \caption{Rúbrica del Bloc A --- Front Matter}
  \label{tab:bloc-a}
  {\footnotesize
  \setlength{\tabcolsep}{4pt}
  \renewcommand{\arraystretch}{1.5}
  \begin{tabularx}{\textwidth}{>{\raggedright\arraybackslash}p{3cm}YYYY}
    \toprule
    \rowcolor{upcblue!15}
    \textbf{Criteri} & \textbf{Excel·lent (0.5 pt)} & \textbf{Notable (0.375 pt)} & \textbf{Aprovat (0.25 pt)} & \textbf{Suspès (0.125 pt)} \\
    \midrule
    A1 Portada & Títol, autor, data i curs presents i ben formats & Falta un element o té errors menors de format & Falten dos elements o el format és deficient & Portada absent o completament incorrecta \\
    A2 Declaració d'autoria & Declaració signada, llista completa d'eines IA, afirmació de comprensió & Declaració present però falta un element (llista IA o afirmació) & Declaració present però incompleta o genèrica & Declaració absent \\
    A3 Abstract trilingüe & Abstract en català, castellà i anglès, cadascun amb 4--6 paraules clau & Abstract en tres idiomes però keywords incompletes o faltants en un idioma & Abstract en menys de tres idiomes o sense keywords & Abstract absent o en un sol idioma sense keywords \\
    A4 TOC + TOT + TOF & Taula de continguts, de taules i de figures presents i actualitzades & Dues de les tres llistes presents & Només una llista present & Cap llista present \\
    \bottomrule
  \end{tabularx}
  }
\end{table}

\section{Bloc B --- Estructura Científica (5 punts)}

La Taula~\ref{tab:bloc-b} defineix la qualitat esperada de l'estructura i del contingut científic del document.

\begin{table}[htbp]
  \centering
  \caption{Rúbrica del Bloc B --- Estructura Científica}
  \label{tab:bloc-b}
  {\footnotesize
  \setlength{\tabcolsep}{4pt}
  \renewcommand{\arraystretch}{1.5}
  \begin{tabularx}{\textwidth}{>{\raggedright\arraybackslash}p{3cm}YYYY}
    \toprule
    \rowcolor{upcblue!15}
    \textbf{Criteri} & \textbf{Excel·lent (100\%)} & \textbf{Notable (75\%)} & \textbf{Aprovat (50\%)} & \textbf{Suspès (25\%)} \\
    \midrule
    B1 Introducció \emph{(0.75 pt)} & Motivació, objectius i estructura del document clarament exposats en capítol propi & Dos dels tres elements presents & Un dels tres elements o capítol molt breu & Introducció absent o no té contingut propi \\
    B2 Capítols del cos \emph{(1.0 pt)} & Mínim 2 capítols entre Introducció i Conclusions, cada un amb mínim 2 seccions i subseccions & 2 capítols presents però seccions o subseccions insuficients & Un sol capítol del cos o estructura molt plana & No hi ha capítols del cos diferenciats \\
    B3 Conclusions crítiques \emph{(1.25 pt)} & Inclou limitacions del tema, preguntes obertes i reflexió personal; va clarament més enllà del resum & Dos dels tres elements crítics presents & Conclusions presents però bàsicament un resum dels capítols anteriors & Conclusions absents o purament descriptives sense cap element crític \\
    B4 Elements de contingut \emph{(1.0 pt)} & Mínim 2 taules i 2 figures, totes referenciades al text amb \texttt{\textbackslash ref\{\}}, totes les cites inline presents & Taules i figures presents però alguna sense cross-reference o cita inline faltant & Taules o figures per sota del mínim, o moltes sense referenciar & Taules i figures absents o cap cross-reference \\
    B5 Referències \emph{(0.5 pt)} & Mínim 8 fonts, totes citades explícitament al text & 8 fonts presents però alguna no citada al text & Menys de 8 fonts o múltiples sense citar & Menys de 4 fonts o bibliografia absent \\
    B6 Apèndix \emph{(0.5 pt)} & Mínim 1 apèndix amb contingut substantiu, referenciat des del text principal & Apèndix present però no referenciat o contingut mínim & Apèndix present però buit o irrellevant & Apèndix absent \\
    \bottomrule
  \end{tabularx}
  }
\end{table}

\section{Bloc C --- Execució LaTeX (3 punts)}

La Taula~\ref{tab:bloc-c} avalua l'ús correcte de LaTeX i la qualitat formal del fitxer lliurat.

\begin{table}[htbp]
  \centering
  \caption{Rúbrica del Bloc C --- Execució LaTeX}
  \label{tab:bloc-c}
  {\footnotesize
  \setlength{\tabcolsep}{4pt}
  \renewcommand{\arraystretch}{1.5}
  \begin{tabularx}{\textwidth}{>{\raggedright\arraybackslash}p{3cm}YYYY}
    \toprule
    \rowcolor{upcblue!15}
    \textbf{Criteri} & \textbf{Excel·lent (100\%)} & \textbf{Notable (75\%)} & \textbf{Aprovat (50\%)} & \textbf{Suspès (25\%)} \\
    \midrule
    C1 Paquets requerits \emph{(1.0 pt)} & Tots els paquets requerits presents i usats correctament: babel (catalan), biblatex o natbib, graphicx, booktabs, geometry, fancyhdr & Tots presents però algun usat incorrectament o amb warnings & Falten 1--2 paquets requerits & Falten 3 o més paquets requerits \\
    C2 Cross-references \emph{(1.0 pt)} & Totes les taules i figures tenen \texttt{\textbackslash label\{\}} i estan referenciades amb \texttt{\textbackslash ref\{\}} al text & La majoria referenciades però alguna omissió & Alguns \texttt{\textbackslash label} i \texttt{\textbackslash ref} presents però ús inconsistent & Cap cross-reference o \texttt{\textbackslash label} usats \\
    C3 Format general \emph{(1.0 pt)} & Marges A4 estàndard (geometry), capçaleres i peus de pàgina configurats (fancyhdr), interlineat adequat & Format majoritàriament correcte amb algun element de fancyhdr incorrecte & Format bàsic present però marges o capçaleres incorrectes & Document sense format: marges per defecte, sense capçaleres \\
    \bottomrule
  \end{tabularx}
  }
\end{table}

\section{Punts Extres --- Ambició i qualitat excepcional}

El professor pot atorgar manualment entre 0.5 i 1.0 punts addicionals a treballs que superin clarament els mínims establerts. Alguns exemples són la inclusió d'un glossari o d'una llista d'acrònims, diagrames originals fets amb TikZ, taules avançades amb \texttt{longtable} o \texttt{multirow}, un estil tipogràfic especialment cuidat o un contingut d'una profunditat o originalitat destacables. Aquests punts extra no estan garantits i depenen exclusivament del criteri del professor.

\section{Càlcul de la nota i lliurament}

La nota del lliurable es calcula amb la fórmula següent:
\[
  \mathrm{Nota\ Lliurable\ 1} =
  \left(\frac{\mathrm{Bloc\ A}}{2}\right)\times 2 +
  \left(\frac{\mathrm{Bloc\ B}}{5}\right)\times 5 +
  \left(\frac{\mathrm{Bloc\ C}}{3}\right)\times 3 +
  \mathrm{Punts\ extres}
\]

La Taula~\ref{tab:resum-punts} resumeix el pes màxim de cada bloc.

\begin{table}[htbp]
  \centering
  \caption{Resum de puntuacions del Lliurable 1}
  \label{tab:resum-punts}
  {\small
  \begin{tabular}{>{\raggedright\arraybackslash}p{10cm}c}
    \toprule
    \rowcolor{upcblue!15}
    \textbf{Bloc} & \textbf{Punts màxims} \\
    \midrule
    Bloc A --- Front Matter & 2.0 \\
    Bloc B --- Estructura Científica & 5.0 \\
    Bloc C --- Execució LaTeX & 3.0 \\
    Total sense bonus & 10.0 \\
    Punts extres (opcional) & fins a +1.0 \\
    \bottomrule
  \end{tabular}
  }
\end{table}

\begin{center}
  \fbox{\parbox{0.94\textwidth}{
    \centering
    Data límit de lliurament: \textbf{\DeadlineDate}\\[0.2cm]
    Lliurar el fitxer \texttt{.pdf} i el codi font \texttt{.tex} (o arxiu \texttt{.zip} amb tots els fitxers) a l'aula virtual.\\[0.2cm]
    Retards no justificats descomptaran 1 punt per dia.
  }}
\end{center}

\end{document}
